\chapter{KẾT LUẬN VÀ HƯỚNG PHÁT TRIỂN}

\section{Kết luận}
Đề tài đã xây dựng hướng tiếp cận phát hiện dấu hiệu sai phạm trong quảng cáo thực phẩm chức năng trên mạng xã hội theo bài toán phân loại đa nhãn. Quy trình gồm thu thập dữ liệu, gán nhãn, huấn luyện mô hình, điều chỉnh ngưỡng và đánh giá kết quả trên từng nhóm dấu hiệu.

\section{Ưu điểm}
\begin{itemize}
    \item Bộ nhãn bám sát các dạng vi phạm thường gặp trong quảng cáo thực phẩm chức năng.
    \item Hướng tiếp cận đa nhãn phù hợp với thực tế một quảng cáo có thể chứa nhiều dấu hiệu sai phạm.
    \item Quy trình có thể mở rộng sang dữ liệu hình ảnh, video và âm thanh.
\end{itemize}

\section{Hạn chế}
\begin{itemize}
    \item Dữ liệu mạng xã hội có nhiều cách diễn đạt lách luật, gây khó khăn cho mô hình.
    \item Một số nhãn có thể bị mất cân bằng dữ liệu, làm giảm F1 macro.
    \item Giai đoạn hiện tại chủ yếu tập trung vào văn bản, chưa khai thác đầy đủ tín hiệu đa phương thức.
\end{itemize}

\section{Hướng phát triển}
\begin{itemize}
    \item Mở rộng bộ dữ liệu bằng cách thu thập thêm mẫu từ nhiều nền tảng và nhóm sản phẩm khác nhau.
    \item Bổ sung OCR để phát hiện chữ nhúng trong ảnh hoặc video.
    \item Bổ sung ASR để phân tích lời thoại trong video quảng cáo.
    \item Kết hợp mô hình nhận diện hình ảnh nhằm phát hiện áo blouse, logo giả mạo hoặc bối cảnh y tế.
    \item Xây dựng giao diện hỗ trợ kiểm duyệt viên xem kết quả, lý do cảnh báo và bằng chứng liên quan.
\end{itemize}
